\section{Critical-layer geometry and a history-selection obstruction}
\label{sec:supplementary-outer-skeleton}

\subsection{Proofs for the preparation-independent critical layer}

\begin{proof}[Proof of \cref{prop:unprepared-outer-skeleton}]
Put \(a_*=D\gamma\).  The estimates
\cref{eq:poisson-bound,eq:shared-products} imply, with constants independent
of \(N\),
\begin{align}
 \|\mathcal N_N(x)\|_N
 &\le C_2\|x\|_N^2+C_3\|x\|_N^3,\notag\\
 \|D\mathcal N_N(x)\|
 &\le C_4\|x\|_N+C_5\|x\|_N^2.                     \label{eq:outer-N-bounds}
\end{align}
For fixed \(r\), consider on \(E_N\)
\begin{equation}
 \mathcal T_{N,r}(z)=A_N^{-1}P_{\perp,N}
       \mathcal N_N(r\mathbf1+z).                   \label{eq:outer-critical-contraction}
\end{equation}
Choose \(M_z\) larger than twice the common quadratic bound in
\cref{eq:outer-N-bounds}.  If \(r_0\) is sufficiently small, uniformly in
\(N\), the map \(\mathcal T_{N,r}\) sends
\(\{\|z\|_N\le M_zr^2\}\) into itself and has Lipschitz constant at most
one half.  Its fixed point is the unique fixed point in that ball and is the
solution \(z_{0,N}(r)\) of \cref{eq:outer-critical-transverse}.

On this ball,
\begin{equation}
 A_N-P_{\perp,N}D\mathcal N_N(r\mathbf1+z_{0,N}(r))
                                                               \label{eq:outer-critical-derivative}
\end{equation}
has inverse norm at most \(2a_*^{-1}\).  Because \(\mathcal N_N\) is a
polynomial with uniform multilinear bounds, induction in the differentiated
fixed-point equation gives, for every fixed finite \(s\), uniform \(C^s\)
bounds.  This proves \cref{eq:outer-critical-high-regularity}; in particular
the \(C^4\) bounds are uniform.  Substituting
\(z_{0,N}(r)=z_{2,N}r^2+z_{3,N}r^3+O_{\rm unif}(r^4)\) into
\cref{eq:outer-critical-transverse} gives
\begin{align*}
 A_Nz_{2,N}&=P_{\perp,N}c_N,\\
 A_Nz_{3,N}&=P_{\perp,N}
       \left(2c_N\circ g_N+\frac\beta3\mathbf1\right).
\end{align*}
Since \(P_{\perp,N}\mathbf1=0\), these identities give
\(z_{2,N}=g_N\) and \cref{eq:outer-g3}.  Taylor's formula applied after
the uniform \(C^4\) estimate yields the first line of
\cref{eq:outer-critical-expansion}.  Stationary projection of
\(\mathcal N_N(r\mathbf1+z_{0,N}(r))\) gives
\begin{equation*}
 \alpha r^2+\left\{2\pi_N^T(c_N\circ g_N)+\frac\beta3\right\}r^3
 +O_{\rm unif}(r^4),
\end{equation*}
which proves the second line and identifies \(\kappa_N\).

The voltage field at the constant history \(V_{0,N}(r)\) is
\begin{equation*}
 (2/3-w_{0,N}(r))\mathbf1
 -\mathcal N_N(x_{0,N}(r))+A_Nz_{0,N}(r).
\end{equation*}
Its stationary and transverse projections vanish by
\cref{eq:outer-critical-transverse,eq:outer-critical-curve}.  Moreover,
every delayed difference is exactly
\(V_{0,N}(r)-V_{0,N}(r)=0\).  This proves the preparation- and
redistribution-free assertion for every \(\delta\).  The recovery law is
irrelevant to this frozen voltage equation.

We next construct the layer splitting without assuming that \(P_N\) is
normal or diagonalizable.  Write
\begin{equation}
 d_N(r)=2c_N\circ x_{0,N}(r)
          +\beta x_{0,N}(r)^{\circ2}.                \label{eq:outer-diagonal-coefficient}
\end{equation}
The decomposition
\(\R^N=\operatorname{span}\{\mathbf1\}\oplus E_N\),
\(x=u\mathbf1+h\), is isometric in the sense that
\(\|x\|_N=|u|+\|h\|_N\).  In these coordinates
\begin{equation}
 \mathcal J^{\rm fv}_{N}(r)=
 \begin{pmatrix}a(r)&b(r)\\c(r)&S(r)\end{pmatrix},             \label{eq:outer-J-blocks}
\end{equation}
where
\begin{align}
 a&=-\pi_N^Td_N,&
 bh&=-\pi_N^T(d_N\circ h),\notag\\
 c&=-P_{\perp,N}d_N,&
 S&=A_N-P_{\perp,N}\diag(d_N)|_{E_N}.               \label{eq:outer-block-formulas}
\end{align}
Because \(x_{0,N}(r)=r\mathbf1+O_{\rm unif}(r^2)\),
\begin{equation}
 |a|+\|b\|+\|c\|+\|S-A_N\|\le C_J|r|              \label{eq:outer-block-smallness}
\end{equation}
with one \(C_J\).  Decrease \(r_0\) until
\(C_Jr_0\le a_*/64\).  Bounded-perturbation variation of constants then
gives
\begin{equation}
 \|e^{S(r)t}\|\le e^{-(a_*-C_J|r|)t},\qquad t\ge0. \label{eq:outer-S-semigroup}
\end{equation}

Seek the center and stable invariant graphs in the forms \(h=H(r)u\) and
\(u=L(r)h\).  Their equations are
\begin{align}
 c+SH&=H(a+bH),\notag\\
 aL+b&=L(cL+S).                                      \label{eq:outer-block-graph-equations}
\end{align}
The linear Sylvester operators at the origin are
\(H\mapsto A_NH\) and \(L\mapsto LA_N\).  Their inverse norms are at most
\(a_*^{-1}\), by the forward semigroup integral and its dual.  More
explicitly, \cref{eq:outer-block-graph-equations} is equivalent to
\begin{align}
 H&=A_N^{-1}\{-c-(S-A_N)H+H(a+bH)\},\notag\\
 L&=\{aL+b-LcL-L(S-A_N)\}A_N^{-1}.                  \label{eq:outer-HL-contractions}
\end{align}
On the balls \(\|H\|,\|L\|\le4C_J|r|/a_*\),
\cref{eq:outer-block-smallness} makes both right-hand sides self-maps with
Lipschitz constant at most one half after one further decrease of \(r_0\).
Hence \(H,L=O_{\rm unif}(r)\), they have uniform finite derivatives, and
\(1-LH\) stays uniformly away from zero.

Set
\begin{equation}
 e_{c,N}(r)=\frac{\mathbf1+H(r)}{1-L(r)H(r)},
 \qquad
 \ell_{c,N}(r)(u\mathbf1+h)=u-L(r)h.                \label{eq:outer-eigen-from-graphs}
\end{equation}
The graph equations show that these are a right and left eigenobject for the
same real scalar eigenvalue \(\lambda_{c,N}=a+bH\), and the two
normalizations in \cref{eq:outer-eigen-normalization} are automatic.  The
stable graph is \(\Ker\ell_{c,N}\), and its generator in the \(h\)
coordinate is \(S+cL\).  In fact,
\begin{equation}
 \|S+cL-A_N\|\le C|r|.                               \label{eq:outer-stable-generator-perturbation}
\end{equation}
The graph coordinate \(h\mapsto Lh\,\mathbf1+h\), its inverse on the stable
graph, and the projection \(P_N^s\) all have norm at most one fixed constant,
because \(L,H=O_{\rm unif}(r)\) and \(1-LH\) is uniformly separated from zero.
Bounded-perturbation variation of constants applied to
\cref{eq:outer-stable-generator-perturbation} therefore gives
\cref{eq:outer-stable-layer-semigroup} with a uniform prefactor.

The first-order parts of \cref{eq:outer-block-graph-equations} give
\begin{equation}
 H(r)=2r g_N+O_{\rm unif}(r^2),
 \qquad
 L(r)=-2r\,\pi_N^T\diag(c_N)A_N^{-1}P_{\perp,N}
          +O_{\rm unif}(r^2).                        \label{eq:outer-HL-expansion}
\end{equation}
Together with \cref{eq:outer-eigen-from-graphs}, these identities prove the
right and left expansions in \cref{eq:outer-eigen-expansion}.

Differentiate the exact critical equation in \(r\).  Since the derivative
with respect to voltage is \(\mathcal J^{\rm fv}_{N}(r)\), one obtains
\begin{equation}
 \mathcal J^{\rm fv}_{N}(r)V_{0,N}'(r)
 =w_{0,N}'(r)\mathbf1.                               \label{eq:outer-critical-differentiated}
\end{equation}
Left multiplication by \(\ell_{c,N}(r)\) gives the exact identity in
\cref{eq:outer-eigen-identity}.  Moreover,
\begin{equation*}
 \Theta_N(r)=1-L(r)z_{0,N}'(r)=1+O_{\rm unif}(r^2).
\end{equation*}
The expansion of \(w_{0,N}'/\Theta_N\) yields the eigenvalue expansion.
It also shows, after decreasing \(r_0\), that the collective eigenvalue is
negative for \(r>0\) and positive for \(r<0\).  The stable complement has
no eigenvalue with nonnegative real part by
\cref{eq:outer-stable-layer-semigroup}, proving the index statements.
\end{proof}

\begin{proof}[Proof of \cref{cor:outer-reduced-actions}]
For sensed
recovery, substitution of
\(x_{0,N}(r)=r\mathbf1+g_Nr^2+O_{\rm unif}(r^3)\) gives
\begin{equation*}
 \sigma_N^{\rm rec}(r)
 =r+\pi_N^T(\varpi_N\circ g_N+d_N^{\rm rec})r^2+O_{\rm unif}(r^3),
\end{equation*}
which is \cref{eq:outer-recovery-expansion}.  In slow time, the recovery
equation restricted formally to the critical curve reads
\begin{equation*}
 w_{0,N}'(r)r_T=\sigma_N^\diamond(r)-\delta^2\nu,
\end{equation*}
and hence gives \cref{eq:outer-reduced-speed}.

Uniformly for the two laws,
\begin{equation}
 w_{0,N}'(r)=-2\alpha r+O_{\rm unif}(r^2),\qquad
 \sigma_N^\diamond(r)=r+O_{\rm unif}(r^2).         \label{eq:outer-speed-inputs}
\end{equation}
First choose \(r_{\rm out}\) so small that, on
\(0<|r|\le r_{\rm out}\),
\begin{equation*}
 \alpha|r|\le|w_{0,N}'(r)|\le3\alpha|r|,
 \qquad
 |\sigma_N^\diamond(r)-r|\le|r|/4.
\end{equation*}
Then decrease \(\delta_0\) so that
\(\delta^2\nu_*\le\rho_\delta/4\) and
\(\rho_\delta\le r_{\rm out}/2\).  The numerator in
\cref{eq:outer-reduced-speed} has the sign of \(r\) and modulus between
\(|r|/2\) and \(3|r|/2\), while \(w_{0,N}'\) has the opposite sign.
This proves \cref{eq:outer-speed-sign} and the displayed expansion.

Combining the exact identities
\cref{eq:outer-eigen-identity,eq:outer-reduced-speed} gives
\begin{equation}
 \frac{\lambda_{c,N}(r)}{q^\diamond_{N,\delta,\nu}(r)}
 =\frac{w_{0,N}'(r)^2}
  {\Theta_N(r)\{\sigma_N^\diamond(r)-\delta^2\nu\}}.          \label{eq:outer-action-identity}
\end{equation}
It has the sign of \(r\), and the preceding bounds imply
\begin{equation}
 \frac{4\alpha^2}{9}|r|
 \le\left|\frac{\lambda_{c,N}(r)}
                 {q^\diamond_{N,\delta,\nu}(r)}\right|
 \le36\alpha^2|r|.                                  \label{eq:outer-action-integrand-bounds}
\end{equation}
Integration proves positivity and
\cref{eq:outer-action-lower-bound}.  Finally,
\cref{eq:outer-speed-inputs,eq:outer-action-identity} gives
\begin{equation*}
 \frac{\lambda_{c,N}(r)}{q^\diamond_{N,\delta,\nu}(r)}
 =4\alpha^2r+O_{\rm unif}(r^2+\delta^2).
\end{equation*}
Integrating on either branch proves
\cref{eq:outer-action-expansion}.  Substituting
\(r=\delta X\), \(|X|=S_\delta/(2\alpha)\), into
\cref{eq:outer-critical-expansion} gives
\cref{eq:outer-critical-fold-offset}.
\end{proof}

\subsection{Positive-\texorpdfstring{\(\delta\)}{delta} frozen-voltage histories}

\begin{proof}[Proof of \cref{prop:frozen-voltage-history-splitting}]
We give a direct trajectory-space proof.  In particular, no determinant,
normality of \(P_N\), or backward RFDE semiflow is used.  The common layer
bound and the structural-ball radius give
\begin{equation}
 b_*:=|K|\sup_{N,\,\|\eta\|_{\mathfrak T_N}<\eta_0}
       \sum_{k=0}^m\|L_{k,N}(\eta)\|<\infty.         \label{eq:frozen-layer-common-bound}
\end{equation}

We first collect the current-state estimates.  Decrease \(r_{\rm out}\) so
that, for one \(c_0>0\),
\begin{equation}
 |\lambda_{c,N}(r)|\ge c_0|r|,
 \qquad 0<|r|\le r_{\rm out}.                       \label{eq:frozen-collective-gap}
\end{equation}
For \(r>0\), the rank-one eigenprojection and
\cref{eq:outer-stable-layer-semigroup} give
\begin{equation}
 \|e^{\delta^{-1}\mathcal J_N^{\rm fv}(r)s}\|
 \le M_0e^{-\mathfrak a_N(r)s/\delta},\qquad s\ge0. \label{eq:frozen-current-attracting}
\end{equation}
For \(r<0\), put
\begin{equation}
 P_N^u(r)=e_{c,N}(r)\ell_{c,N}(r),\qquad
 P_N^s(r)=\Id-P_N^u(r).                              \label{eq:frozen-current-projections}
\end{equation}
The same estimates imply
\begin{align}
 \|e^{\delta^{-1}\mathcal J_N^{\rm fv}(r)s}P_N^s(r)\|
 &\le M_0e^{-\mathfrak a_N(r)s/\delta},&&s\ge0,\notag\\
 \|e^{\delta^{-1}\mathcal J_N^{\rm fv}(r)s}P_N^u(r)\|
 &\le M_0e^{\mathfrak a_N(r)s/\delta},&&s\le0.     \label{eq:frozen-current-repelling}
\end{align}
The harmonic definition in \cref{eq:frozen-history-a} is only a smooth
common lower rate for the two current blocks.  Uniform boundedness of the
projections follows from \cref{eq:outer-eigen-from-graphs}.

Write \(\kappa=\kappa_{N,\delta}(r)\) for the rest of the proof.  From
\cref{eq:frozen-history-rate},
\begin{equation}
 2\kappa<\frac{\mathfrak a_N(r)}{2\delta},\qquad
 e^{2\kappa\theta_m}\le\delta^{-2\vartheta}.        \label{eq:frozen-cap-bounds}
\end{equation}
If a current path \(y\) and a prescribed old history are assembled into
\(\widehat y\), set
\begin{equation}
 \mathcal R_\eta[\widehat y_s]
 =\delta K\sum_{k=0}^mL_{k,N}(\eta)
       \{\widehat y(s)-\widehat y(s-\theta_k)\}.     \label{eq:frozen-delay-operator}
\end{equation}
On any exponentially weighted trajectory space, translation across one
delay gives at most the factor \(e^{2\kappa\theta_m}\).  The current Green
integrals and \cref{eq:frozen-cap-bounds} therefore give the common
perturbation number
\begin{equation}
 q_{N,\delta}(r)
 \le C\frac{\delta^2
       \{1+e^{2\kappa\theta_m}\}}{\mathfrak a_N(r)}.
                                                               \label{eq:frozen-LP-number}
\end{equation}
Since \(\mathfrak a_N(r)\ge c|r|\) on the retained interval and
\(|r|\ge\rho_\delta=\delta S_\delta/(2\alpha)\),
\begin{equation}
 q_{N,\delta}(r)
 \le C\frac{\delta^{1-2\vartheta}}{S_\delta}=o(1).  \label{eq:frozen-LP-small}
\end{equation}
This proves \cref{eq:frozen-history-smallness} and permits one common choice
of \(\delta_0\) for which every contraction below has constant at most
one half.

Suppose first that \(r>0\).  For \(\phi\in\mathscr X_N\), prescribe
\(\widehat y(\theta)=\phi(\theta)\) for \(-\theta_m\le\theta\le0\) and use
variation of constants for \(s\ge0\):
\begin{equation}
 y(s)=e^{\delta^{-1}\mathcal J_N^{\rm fv}(r)s}\phi(0)
 +\int_0^s e^{\delta^{-1}\mathcal J_N^{\rm fv}(r)(s-u)}
       \mathcal R_\eta[\widehat y_u]\,\dd u.        \label{eq:frozen-attracting-VoC}
\end{equation}
For a future path put
\(\|y\|_+=\sup_{s\ge0}e^{\kappa s}\|y(s)\|_N\).
If \(u-\theta_k\ge0\), then
\begin{equation*}
 e^{\kappa u}\|y(u-\theta_k)\|_N
 \le e^{\kappa\theta_m}\|y\|_+.
\end{equation*}
If \(u-\theta_k<0\), the factor two in
\cref{eq:frozen-history-norm} gives instead
\begin{equation}
 e^{\kappa u}\|\phi(u-\theta_k)\|_N
 \le e^{-\kappa u+2\kappa\theta_k}
       \|\phi\|_{N,r,\delta}^{\sharp}
 \le e^{2\kappa\theta_m}
       \|\phi\|_{N,r,\delta}^{\sharp}.             \label{eq:frozen-old-history-bound}
\end{equation}
Using \cref{eq:frozen-current-attracting} in
\cref{eq:frozen-attracting-VoC}, the Green integral is bounded by
\begin{equation*}
 \int_0^s e^{-\mathfrak a_N(r)(s-u)/\delta}
              e^{-\kappa u}\,\dd u
 \le\frac{e^{-\kappa s}}
          {\mathfrak a_N(r)/\delta-\kappa}.
\end{equation*}
Thus \cref{eq:frozen-LP-number} and a Neumann argument give
\begin{equation}
 \|y(s)\|_N
 \le C e^{-\kappa s}\|\phi\|_{N,r,\delta}^{\sharp}.
                                                               \label{eq:frozen-attracting-current-decay}
\end{equation}
To lift this estimate to a history segment, take
\(-\theta_m\le\theta\le0\).  If \(s+\theta\ge0\), then
\begin{equation*}
 e^{2\kappa\theta}\|y(s+\theta)\|_N
 \le Ce^{-\kappa s}e^{\kappa\theta}
       \|\phi\|_{N,r,\delta}^{\sharp}.
\end{equation*}
If \(s+\theta<0\), translation of the prescribed history gives
\begin{equation*}
 e^{2\kappa\theta}\|\phi(s+\theta)\|_N
 \le e^{-2\kappa s}\|\phi\|_{N,r,\delta}^{\sharp}.
\end{equation*}
Taking the supremum proves
\cref{eq:frozen-attracting-history-estimate}.  Notice that the old-history
shift is estimated directly; a spectral bound alone would not prove this
semigroup estimate.

Now let \(r<0\).  Define the unperturbed unstable history embedding
\begin{equation}
 (\jmath_N^ua)(\theta)
 =e^{\delta^{-1}\mathcal J_N^{\rm fv}(r)\theta}a,\qquad
 a\in\Ran P_N^u(r),\quad -\theta_m\le\theta\le0,    \label{eq:frozen-base-unstable-embedding}
\end{equation}
and the fixed complementary coordinate space
\begin{equation}
 \mathscr K_N(r)=\{\psi\in\mathscr X_N:
                         P_N^u(r)\psi(0)=0\}.       \label{eq:frozen-base-stable-coordinate}
\end{equation}
The map \(\jmath_N^u\), current evaluation, and both coordinate
projections are bounded uniformly in the fading norm.  Every history has a
unique representation \(\psi+\jmath_N^ua\).

For a given \(\psi\in\mathscr K_N(r)\), seek
\(a\in\Ran P_N^u(r)\) and a future path \(y\) with \(\|y\|_+<\infty\).
Assemble the old history as
\(\widehat y(\theta)=\psi(\theta)+(\jmath_N^ua)(\theta)\).  The
Lyapunov--Perron equations are
\begin{align}
 y(s)={}&e^{\delta^{-1}\mathcal J_N^{\rm fv}(r)s}
              P_N^s(r)\psi(0)\notag\\
 &+\int_0^s e^{\delta^{-1}\mathcal J_N^{\rm fv}(r)(s-u)}P_N^s(r)
       \mathcal R_\eta[\widehat y_u]\,\dd u\notag\\
 &-\int_s^\infty e^{\delta^{-1}\mathcal J_N^{\rm fv}(r)(s-u)}P_N^u(r)
       \mathcal R_\eta[\widehat y_u]\,\dd u,        \label{eq:frozen-stable-LP}\\
 a={}&-\int_0^\infty
       e^{-\delta^{-1}\mathcal J_N^{\rm fv}(r)u}P_N^u(r)
       \mathcal R_\eta[\widehat y_u]\,\dd u.        \label{eq:frozen-stable-boundary}
\end{align}
The last identity is simply the unstable projection of the first one at
\(s=0\), so endpoint compatibility is exact.  The stable forward integral
has denominator
\(\mathfrak a_N(r)/\delta-\kappa\), and the unstable future integral has
denominator \(\mathfrak a_N(r)/\delta+\kappa\).  Together with
\cref{eq:frozen-old-history-bound}, this makes
\cref{eq:frozen-stable-LP,eq:frozen-stable-boundary} a contraction with
constant bounded by \cref{eq:frozen-LP-number}.  Its unique solution obeys
\begin{equation}
 \|y\|_++\|a\|_N
 \le C\|\psi\|_{N,r,\delta}^{\sharp}.               \label{eq:frozen-stable-LP-bound}
\end{equation}
Linearity defines a bounded graph map
\(F^s_{N,\delta,r,\eta}:\mathscr K_N(r)\to\Ran P_N^u(r)\), with
\begin{equation}
 \|F^s_{N,\delta,r,\eta}\|
 \le Cq_{N,\delta}(r),                               \label{eq:frozen-stable-graph-small}
\end{equation}
and
\begin{equation}
 \mathscr W^s_{N,\delta,r,\eta}
 =\{\psi+\jmath_N^uF^s_{N,\delta,r,\eta}\psi:
                 \psi\in\mathscr K_N(r)\}.         \label{eq:frozen-stable-history-graph}
\end{equation}
Conversely, variation of constants for any solution with finite
\(\|y\|_+\), followed by unstable projection and passage of the terminal
time to infinity, recovers
\cref{eq:frozen-stable-LP,eq:frozen-stable-boundary}.  The graph therefore
characterizes exactly the histories whose future decays at this rate.  It is
invariant by uniqueness.  The same two-case history-segment estimate used
above proves \cref{eq:frozen-repelling-stable-estimate}.

For the complementary line, solve directly on \((-\infty,0]\).  Given
\(a\in\Ran P_N^u(r)\), seek \(y\) satisfying
\begin{align}
 y(s)={}&e^{\delta^{-1}\mathcal J_N^{\rm fv}(r)s}a\notag\\
 &+\int_0^s e^{\delta^{-1}\mathcal J_N^{\rm fv}(r)(s-u)}P_N^u(r)
       \mathcal R_\eta[y_u]\,\dd u\notag\\
 &+\int_{-\infty}^s
       e^{\delta^{-1}\mathcal J_N^{\rm fv}(r)(s-u)}P_N^s(r)
       \mathcal R_\eta[y_u]\,\dd u.                 \label{eq:frozen-unstable-LP}
\end{align}
On the space with norm
\(\sup_{s\le0}e^{-\kappa s}\|y(s)\|_N\), a delayed value lies farther in
the decaying direction:
\begin{equation*}
 e^{-\kappa s}\|y(s-\theta_k)\|_N
 \le e^{-\kappa\theta_k}
       \sup_{u\le0}e^{-\kappa u}\|y(u)\|_N.
\end{equation*}
The two Green integrals hence give a contraction no larger than
\cref{eq:frozen-LP-number}.  This constructs a one-dimensional space of
complete backward histories, which is invariant by autonomy and uniqueness.
For its history at zero, the same contraction estimate gives explicitly
\begin{equation}
 \|(y_a)_0-\jmath_N^ua\|_{N,r,\delta}^{\sharp}
 \le Cq_{N,\delta}(r)\|a\|_N.                       \label{eq:frozen-unstable-graph-small}
\end{equation}
The stable and unstable graphs are \(o(1)\) perturbations of complementary
base graphs, so they remain complementary and their projections are bounded
by one common constant.  This proves
\cref{eq:frozen-repelling-history-splitting}.

For completeness, invariance of the one-dimensional line gives a continuous
one-dimensional group.  Hence there is a real \(\lambda^u\) and a normalized
current vector \(e^u\), with \(\ell_{c,N}(r)e^u=1\), for which its histories
have the form \cref{eq:frozen-unstable-history}.  Projecting
\cref{eq:frozen-unstable-LP} onto the current stable graph gives
\begin{equation}
 \|e^u-e_{c,N}(r)\|_N\le C\delta^2.                 \label{eq:frozen-unstable-vector-proof-bound}
\end{equation}
Indeed, the stable current Green denominator is
\(\lambda^u+D\gamma/(2\delta)\), while the delayed coefficient is
\(O(\delta)\); hence the stable resolvent is \(O(\delta)\), without first
using the root expansion.  Substitution of
\(e^{\lambda^us}e^u\) into \cref{eq:frozen-voltage-rfde}, followed by
\(\ell_{c,N}(r)\), yields
\begin{equation}
 \lambda^u-\delta^{-1}\lambda_{c,N}(r)
 =\delta K\sum_{k=0}^m\ell_{c,N}(r)L_{k,N}(\eta)
       \{1-e^{-\lambda^u\theta_k}\}e^u.            \label{eq:frozen-unstable-scalar-equation}
\end{equation}
The right-hand side is \(O(\delta)\), because \(\lambda^u>0\) for small
\(\delta\).  This proves \cref{eq:frozen-unstable-root-estimate}.  The group
on the line now gives
\cref{eq:frozen-unstable-inverse-estimate}.

It remains to justify the structural derivatives.  The delay locations are
fixed, the map \(\eta\mapsto(L_{0,N}(\eta),\ldots,L_{m,N}(\eta))\) is
affine in the full operator-sum norm, and the fading weight is independent of
\(\eta\).  Thus every Lyapunov--Perron map above is twice Fr\'{e}chet
differentiable on one fixed Banach space, and its first two derivatives have
common bounds.  Differentiating the contraction equation once and twice and
using \((1-q_{N,\delta}(r))^{-1}\le2\) gives common bounds for the graph
maps.  The usual graph-projection formulas give the same bounds for the
projections.  Differentiating
\cref{eq:frozen-unstable-scalar-equation} and the stable current equation
gives the asserted bounds for \(\lambda^u,e^u\), and
\cref{eq:frozen-unstable-history}.  More explicitly, solve the stable
current equation for \(e^u=e_{c,N}+e^s(\lambda,\eta)\); its stable resolvent
and its first two \((\lambda,\eta)\)-derivatives are uniformly bounded at the
displayed scales.  After moving the right side of
\cref{eq:frozen-unstable-scalar-equation} to the left, the derivative of the
resulting scalar residual with respect to \(\lambda\) is
\(1+O(\delta)\), because the atoms \(\theta_k\) are fixed.  It is at least
one half for small \(\delta\), so the ordinary implicit-function theorem
supplies the claimed \(C^2_\eta\) root and vector.  No moving evaluation atom
is differentiated.

Finally, complexify the phase space and all graph maps.  Every characteristic
solution is then a phase-space orbit and splits according to
\cref{eq:frozen-repelling-history-splitting}.  The stable estimate excludes a
characteristic exponent with real part greater than \(-\kappa\), except for
the displayed unstable line; the attracting estimate gives the analogous
exclusion there.  A generalized vector at \(\lambda^u\) would produce a
second linearly independent, nondecaying generalized history, contradicting
the one-dimensional unstable summand: indeed,
\cref{eq:frozen-unstable-root-estimate,eq:frozen-history-rate} give
\(\lambda^u>\kappa\), so every polynomial multiple of
\(e^{\lambda^us}\) lies in the backward weighted class used in
\cref{eq:frozen-unstable-LP}.  Thus the displayed real root is
algebraically simple.  This completes the proof.
\end{proof}

\begin{proof}[Proof of \cref{prop:exact-outer-history-equation}]
If \(T=\delta^2t\), then
\[
 \frac{\dd}{\dd t}V(r(T))
 =\delta^2q(r)V'(r),\qquad
 \frac{\dd}{\dd t}w(r(T))
 =\delta^2q(r)w'(r).
\]
The physical delay \(\theta_k/\delta\) in original time is the slow-time
delay \(\delta\theta_k\).  Therefore
\[
 r\!\left(T-\delta\theta_k\right)
 =\Phi_q^{-\delta\theta_k}(r(T))=r_k(r(T)),
\]
and the delayed voltage in the raw RFDE is \(V(r_k(r(T)))\).  Substitute
\(V=\mathbf1+r\mathbf1+z\), use
\(\mathcal A_N\mathbf1=0\), and note that
\(\mathcal A_Nz=A_Nz\).  The voltage equation is then exactly
\cref{eq:exact-outer-voltage-invariance}.  Since
\[
 \pi_N^TV(r)-a
 =1+r+\pi_N^Tz(r)-\{1+\delta^2\nu\}
 =r-\delta^2\nu,
\]
the resource equation is exactly
\cref{eq:exact-outer-resource-invariance}.  Since every retained \(r\) may
serve as the initial value \(r_0\), evaluating the generated equation at
\(t=0\) proves necessity at every retained point.  The same substitution in
reverse proves sufficiency on every orbit segment for which all displayed
points remain in \(J\).

It remains to justify the regularity used below.  From
\cref{eq:exact-outer-resource-invariance},
\[
 w'(r)=\frac{r-\delta^2\nu}{q(r)}
\]
is \(C^1\), so \(w\) is \(C^2\) in the interior.  The scalar flow and each
backtrack \(r_k\) are \(C^1\).  Consequently the right-hand side of
\cref{eq:exact-outer-voltage-invariance} is \(C^1\); division by the nonzero
\(q\) shows that \(z'\), hence \(V'\), is \(C^1\).  Thus \(V\) is \(C^2\)
there as claimed.

For fold time \(s=\delta t\), an offset \(\theta\) corresponds to the
slow-time offset \(\delta\theta\), which proves
\cref{eq:exact-outer-history-embedding}.  Current evaluation and
\(\pi_N^TV'(r)=1+\pi_N^Tz'(r)=1\) give
\cref{eq:exact-outer-history-phase}.  Finally,
\(T=\delta s\).  Defining \(r_*(s)=r(\delta s)\) gives
\(r_*'(s)=\delta q(r_*(s))\), and differentiation of the entire history
embedding gives
\[
 \frac{\dd}{\dd s}\mathcal H_{V,w,q}(r_*(s))
 =\delta q(r_*(s))\partial_r\mathcal H_{V,w,q}(r_*(s)).
\]
Division by its stationary current coordinate proves
\cref{eq:exact-outer-history-tangent}.
\end{proof}

\begin{proof}[Proof of \cref{prop:physical-backtrack-calculus}]
The flow definition is equivalent to
\begin{equation}
 \int_r^{r_j(r)}\frac{\dd u}{q(u)}=-\delta\theta_j.
                                                               \label{eq:backtrack-integral-identity}
\end{equation}
The sign and lower bound in \cref{eq:backtrack-q-class} give the first
estimate in \cref{eq:backtrack-basic-identities}.  Differentiation with
respect to \(r\) gives
\[
 -\frac1{q(r)}+\frac{r_j'(r)}{q(r_j(r))}=0,
\]
which is its second identity.  Differentiation of
\cref{eq:backtrack-integral-identity} in a direction \(h\) gives
\[
 \frac{D_qr_j[h]}{q(r_j)}
 -\int_r^{r_j}\frac{h(u)}{q(u)^2}\,\dd u=0,
\]
and hence \cref{eq:backtrack-first-variation}.

We next prove the differentiability without losing an unnecessary third
spatial derivative of \(q\).  For fixed
\(0\le\tau\le\delta\theta_m\), set
\[
 R_q(\tau,r)=\Phi_q^{-\tau}(r)
\]
and consider the fixed-domain equation
\begin{equation}
 \mathscr H_\tau(q,R)(r)
 =\int_r^{R(r)}\frac{\dd u}{q(u)}+\tau=0.           \label{eq:backtrack-implicit-map}
\end{equation}
The strict inclusion \(J_0\Subset J\) persists under sufficiently small
\(C^1\) perturbations of \(q\) and \(R\).  Hence this equation is defined on
an open product neighborhood.  More explicitly, writing
\(C^a_\rho(K)\) for \(C^a(K)\) with
\cref{eq:backtrack-weighted-norm}, it is a map
\[
 \mathscr H_\tau:
 C^{k+2}_\rho(J)\times\mathcal U_k(I,J_0)longrightarrow C^k_\rho(I),
\]
where \(\mathcal U_k(I,J_0)\) is the open set of \(C^k\) maps whose range
lies in \(\operatorname{int}J_0\).

On either sign component set \(y=\log|r|\).  For each fixed \(k\), the norm
\cref{eq:backtrack-weighted-norm} is uniformly equivalent to
\begin{equation}
 \max_{0\le a\le k}\sup_y
   \|\partial_y^a f(\operatorname{sgn}(J)e^y)\|,    \label{eq:backtrack-log-norm}
\end{equation}
because \((r\partial_r)^a\) is a triangular combination of
\(r^b\partial_r^b\), \(b\le a\).  In these equivalent weighted spaces the
map \(\mathscr H_\tau\) is \(C^3\): three derivatives with respect to its
upper endpoint differentiate \(q^{-1}\) only twice, and therefore
\(q,h\in C^{k+2}\) suffice after \(k\) spatial derivatives.  Moreover,
\[
 D_R\mathscr H_\tau(q,R)v=\frac{v}{q(R)}.
\]
Its inverse is multiplication by \(q(R)\), with norm bounded solely by the
common constants in \cref{eq:backtrack-q-class}.  The Banach implicit-function
theorem therefore makes \(q\mapsto R_q(\tau,\cdot)\) three times
Fr\'{e}chet differentiable in the stated spaces.

Differentiating \cref{eq:backtrack-implicit-map} with respect to \(\tau\)
gives \(\partial_\tau R_q=-q(R_q)\).  Since \(R_q(0,r)=r\), integration in
\(\tau\) proves the first line below.  For the parameter derivatives,
differentiate \cref{eq:backtrack-implicit-map} successively.  Every pure
\(q\)-derivative is an integral over \([r,R_q]\), hence has a factor
\(C\tau\); every endpoint term contains a lower variation of \(R_q\).
Induction on the parameter order, followed by multiplication by the uniformly
bounded inverse of \(D_R\mathscr H_\tau\), gives, for \(1\le b\le3\),
\begin{align}
 \max_{a\le k}\sup_y
 \left|\partial_y^a\{R_q(\tau,r)-r\}\right|
 &\le C\tau,\notag\\
 \max_{a\le k}\sup_y
 \left|\partial_y^aD_q^bR_q(\tau,r)
          [h_1,\ldots,h_b]\right|
 &\le C\tau\prod_{\ell=1}^b
       \|h_\ell\|_{k+2,\rho;J}.                    \label{eq:backtrack-fixed-time-jets}
\end{align}
Here the integral and endpoint Nemytskii estimates are uniform because
\(|R_q/r|\) stays between two fixed positive constants when
\(\epsilon_\rho\) is small.  This also proves directly that the constants do
not depend on \(\rho,\delta\), or the interval locations.  Taking
\(\tau=\delta\theta_j\) and using
\cref{eq:backtrack-log-norm} proves
\cref{eq:backtrack-map-bound,eq:backtrack-frechet-bounds}.

The exact second variation, useful for checking the signs, is
\begin{align}
 D_q^2r_j[h_1,h_2]
={}&\{h_2(r_j)+q'(r_j)R_2\}
       \int_r^{r_j}\frac{h_1}{q^2}
 +\frac{h_1(r_j)}{q(r_j)}R_2\notag\\
 &-2q(r_j)\int_r^{r_j}\frac{h_1h_2}{q^3},
 \qquad R_2=D_qr_j[h_2].                            \label{eq:backtrack-second-variation}
\end{align}
It is the derivative of \cref{eq:backtrack-first-variation}; one further
derivative is covered by \cref{eq:backtrack-fixed-time-jets}.

The collar geometry and the first displacement bound give
\begin{equation}
 \frac12\le\frac{|r_j(r)|}{|r|}\le\frac32,\qquad
 \left|\log\frac{|r_j(r)|}{|r|}\right|
 \le C\epsilon_\rho,                               \label{eq:backtrack-log-shift}
\end{equation}
after decreasing \(\epsilon_0\).  In fact, if
\(Y_q(\tau,y)=\log|R_q(\tau,\operatorname{sgn}(J)e^y)|\), then
\cref{eq:backtrack-fixed-time-jets} and smoothness of the logarithm on the
displayed ratio interval give the stronger bound
\begin{equation}
 \|Y_q(\tau,\cdot)-\Id\|_{C^k_y}
 +\sum_{b=1}^3\|D_q^bY_q(\tau,\cdot)\|_{
      \mathcal L^b(C^{k+2}_\rho(J);C^k_y)}
 \le C\frac{\tau}{\rho},                           \label{eq:backtrack-log-jet-shift}
\end{equation}
Fa\`a di Bruno's formula in the \(y\)-coordinate now proves
\cref{eq:backtrack-composition-bound}.  Apply the fundamental theorem of
calculus to the full \(C^k_y\) displacement in
\cref{eq:backtrack-log-jet-shift} to obtain
\cref{eq:backtrack-composition-difference}.  For every positive parameter
order, each Fa\`a di Bruno term contains at least one variation from
\cref{eq:backtrack-fixed-time-jets}.  A term with \(b\) such factors is
bounded first by \(C\epsilon_\rho^b\); since
\(\epsilon_\rho\le\epsilon_0\), it is at most \(C\epsilon_\rho\).
This proves all three lines of
\cref{eq:backtrack-composition-frechet}, including the third derivative.
The same calculation applied to the joint map \((f,q)\mapsto f\circ r_j(q)\)
has either no variation of \(f\), or exactly one such variation.  Applying it
instead to
\[
 f(\operatorname{sgn}(J)e^{Y_q(y)})
   -f(\operatorname{sgn}(J)e^y)
\]
leaves a factor from \cref{eq:backtrack-log-jet-shift} in every term and
proves \cref{eq:backtrack-joint-difference}; dropping that factor gives the
asserted bound for \(\mathcal C_j\).  The rectangular parameter statement is
now the third-order joint Banach-space chain rule.  Substitution of
\(\rho_\delta=\delta S_\delta/(2\alpha)\) proves
\cref{eq:backtrack-logarithmic-loss}.
\end{proof}

\begin{proof}[Proof of \cref{prop:principal-tracker-endpoint-green}]
The high regularity in
\cref{eq:outer-critical-high-regularity}, the polynomial formula for
\(D\mathcal N_N\), and the critical expansions give
\cref{eq:principal-green-coefficients}.  The sign and zeroth-order speed
bounds are \cref{eq:outer-speed-sign,eq:outer-speed-inputs}.  Repeated
differentiation of
\[
 q_0(r)=\frac{r-\varepsilon\nu}{w_{0,N}'(r)}
\]
and \(|r|\ge\rho_\delta\) gives
\cref{eq:principal-green-q-bounds}; every apparently singular term contains
at most a power of \(\varepsilon/|r|\).  In particular, differentiating
\(q_0\) uses the first, but not the second, of the two ratios
\begin{equation}
 \frac{\varepsilon}{\rho_\delta}
 =\frac{2\alpha\delta}{S_\delta},
 \qquad
 \frac{\varepsilon}{\rho_\delta^2}
 =\frac{4\alpha^2}{S_\delta^2}.                    \label{eq:principal-green-small-ratios}
\end{equation}
After decreasing \(\delta_0\), both ratios are uniformly bounded and tend
to zero.  The second ratio enters later, when the scalar Green equation is
normalized by \(|w_{0,N}'|\asymp|r|\); it is not produced by
differentiating \(q_0\).

Write a dot for differentiation along \(\dot r=q_0(r)\).  The coordinate
\cref{eq:principal-green-p-coordinate} gives the exact system
\begin{align}
 \varepsilon\dot Z
 &=C_N(r)Z-z_{0,N}'(r)\mathfrak p+f,\notag\\
 \varepsilon\dot{\mathfrak p}
 &=w_{0,N}'(r)\mathfrak p
   -w_{0,N}'(r)\ell_N(r)Z-\varepsilon g,\notag\\
 C_N(r)&=B_N(r)+z_{0,N}'(r)\ell_N(r).              \label{eq:principal-green-transformed-system}
\end{align}
Indeed, the second row of
\cref{eq:principal-tracker-operator} reads
\[
 g=-\frac{\dd}{\dd T}\{\varepsilon Q+\ell_NZ\}+w_{0,N}'Q,
\]
which proves the second equation; the first follows from
\(\varepsilon Q=\mathfrak p-\ell_NZ\).  This is the cancellation that the
untransformed derivative coupling conceals.

The transverse Markov estimate gives
\(\|e^{A_NT}\|_{E_N\to E_N}\le e^{-D\gamma T}\) for \(T\ge0\).
Since \(\|C_N(r)-A_N\|\le Cr_{\rm out}\), bounded-perturbation variation of
constants yields the evolution estimate
\begin{equation}
 \|U_C(T,S)\|
 \le C e^{-D\gamma(T-S)/(2\varepsilon)},
 \qquad T\ge S.                                    \label{eq:principal-green-C-evolution}
\end{equation}
The scalar evolution is
\begin{equation}
 u_q(T,S)=\exp\left\{\varepsilon^{-1}
      \int_S^T w_{0,N}'(r(\sigma))\,\dd\sigma\right\}.       \label{eq:principal-green-scalar-evolution}
\end{equation}
It is forward stable on \(I_\delta^a\); on \(I_\delta^r\), its restriction
to the scalar line has a backward-stable inverse.  No ambient backward RFDE
evolution is being used.

We construct one slow particular solution.  Pass to
\(x=\log(|r|/\rho_\delta)\) and first apply a fixed-width linear jet
extension to the right-hand sides and coefficients at each endpoint.  At
the far side of that collar, freeze the coefficients, preserve the stable
matrix bound and the sign of \(w_{0,N}'\), and attach an artificial
\(T\)-halfline that is no longer required to have the form
\(r=\pm\rho_\delta e^x\).  Extend the data with the same linear operator
and freeze the endpoint weight.  Thus the repelling \(Z\)-row has a genuine
past halfline and its scalar row has a genuine future halfline; on the
attracting branch both rows have genuine past halflines.  The extensions
have the same weighted \(C^s\) bounds, uniformly in the length of the
physical logarithmic interval.

Writing tildes for the extended objects, the three diagonal Green kernels
are explicitly
\begin{align}
 (\mathscr G_ZF)(T)
 &=\varepsilon^{-1}\int_{-\infty}^T
       \widetilde U_C(T,S)\widetilde F(S)\,\dd S,\notag\\
 (\mathscr G_{\mathfrak p}^aH)(T)
 &=\varepsilon^{-1}\int_{-\infty}^T
       \widetilde u_q(T,S)\widetilde H(S)\,\dd S,\notag\\
 (\mathscr G_{\mathfrak p}^rH)(T)
 &=-\varepsilon^{-1}\int_T^\infty
       \widetilde u_q(T,S)\widetilde H(S)\,\dd S.  \label{eq:principal-green-diagonal-kernels}
\end{align}
Use \(\mathscr G_Z\) for the transverse row on both branches,
\(\mathscr G_{\mathfrak p}^a\) on the attracting branch, and
\(\mathscr G_{\mathfrak p}^r\) on the repelling branch.  Their block fixed
point has inputs
\(F=f-z_{0,N}'\mathfrak p\) and
\(H=-w_{0,N}'\ell_NZ-\varepsilon g\).
Restriction to the retained interval solves the original equation.  The
halfline is only a device for selecting a bounded right inverse; it is not a
history of the RFDE and is not a physical preparation.

For clarity, the logarithmic norm equivalences behind the estimate are
\begin{align}
 \|Z\|_{\mathsf Z_\varepsilon^s}
 &\asymp\varepsilon^{-1}\|Z\|_{C_x^s},&
 \|\mathfrak p\|_{\mathsf P_\varepsilon^s}
 &\asymp\varepsilon^{-1}\|\mathfrak p/r\|_{C_x^s},&
 \|Q\|_{\mathsf Q_\varepsilon^s}
 &\asymp\varepsilon^{-1}\|rQ\|_{C_x^s}.          \label{eq:principal-green-log-norms}
\end{align}
In these norms the direct scalar source obeys
\begin{equation}
 \|\mathscr G_{\mathfrak p}^{a/r}(-\varepsilon g)\|_{
       \mathsf P_\varepsilon^s}
 \le C_s\frac{\varepsilon}{\rho_\delta^2}
       \|(0,g)\|_{\mathsf Y_\varepsilon^s},       \label{eq:principal-green-g-source}
\end{equation}
because the scalar denominator is \(|w_{0,N}'|\asymp|r|\).  The two
off-diagonal Green maps satisfy
\begin{equation}
 \|\mathcal K_{Z\mathfrak p}\|\le Cr_{\rm out}^2,
 \qquad
 \|\mathcal K_{\mathfrak p Z}\|\le C,
 \qquad
 \|\mathcal K_{Z\mathfrak p}\mathcal K_{\mathfrak p Z}\|
 \le Cr_{\rm out}^2.                               \label{eq:principal-green-loop}
\end{equation}
The first estimate uses \(z_{0,N}'=O(r)\),
\(\mathfrak p=O(\varepsilon|r|)\), and the fixed transverse Green
denominator.  The second uses
\(w_{0,N}'\ell_N=O(r^2)\) and the scalar Green denominator
\(|w_{0,N}'|\asymp|r|\).  Hence the block system is invertible after one
fixed decrease of \(r_{\rm out}\).

To make the derivative induction explicit, put
\(\mathcal D=r\partial_r=(r/q_0)\partial_T\).  Commuting
\(\mathcal D^j\) through the equations produces three kinds of lower-order
terms.  Derivatives of
\(C_N-A_N,z_{0,N}',\ell_N,w_{0,N}'\) retain their respective structural
sizes and hence the same diagonal perturbation and off-diagonal bounds as in
\cref{eq:principal-green-loop}.  Ratios such as
\(\mathcal Dw_{0,N}'/w_{0,N}'\) are merely \(O(1)\); they multiply
lower-order unknowns already controlled by induction and need no smallness.
Only commutators with the transport coefficient
\(\varepsilon q_0/r\) carry \(C_s\varepsilon/|r|\).  After the scalar row
is normalized by both \(|w_{0,N}'|\asymp|r|\) and the
\(\mathfrak p/r\) coordinate, its singular transport ratio is
\(C_s\varepsilon/|r|^2\).  The triangular equivalence between
\(\mathcal D^j\) and \(r^j\partial_r^j\), together with
\cref{eq:principal-green-small-ratios}, lets these terms be absorbed at
each derivative order.  The direct estimates
\cref{eq:principal-green-g-source,eq:principal-green-loop} then close the
same two-by-two block induction successively through order \(s\).  This proves the
\(\mathsf Z_\varepsilon^s\times\mathsf P_\varepsilon^s\) part of
\cref{eq:principal-green-slow-bound}.

To recover the original speed component, do not select another solution of
a scalar equation.  The transformed second row gives the exact identity
\begin{equation}
 Q=\frac{\mathfrak p-\ell_NZ}{\varepsilon}
  =\frac{g+q_0\mathfrak p'}{w_{0,N}'}.             \label{eq:principal-green-Q-recovery}
\end{equation}
If \(\mathfrak p\in\mathsf P_\varepsilon^s\), then the numerator obeys
\begin{equation}
 \left|\partial_r^j\{g+q_0\mathfrak p'\}\right|
 \le C_s\varepsilon|r|^{-j},\qquad0\le j\le s-1. \label{eq:principal-green-Q-forcing}
\end{equation}
The coefficient bounds for \(1/w_{0,N}'\), whose \(j\)-th derivative is
\(O(|r|^{-j-1})\), prove the
\(\mathsf Q_\varepsilon^{s-1}\) estimate.  This is the algebraic component
of the same transformed solution, so it solves both rows of the original
operator with no uniqueness bridge left implicit.  With one additional
input and coefficient derivative, the transformed construction gives
\(\mathfrak p\in\mathsf P_\varepsilon^{s+1}\), and the same identity proves
\cref{eq:principal-green-extra-q-derivative}.

It remains to impose finite endpoint traces.  For the homogeneous transformed
system, the attracting lift is the fixed point of
\begin{align}
 Z(T)&=U_C(T,0)b_z-\varepsilon^{-1}
   \int_0^TU_C(T,S)z_{0,N}'(S)\mathfrak p(S)\,\dd S,\notag\\
 \mathfrak p(T)&=u_q(T,0)b_p-\varepsilon^{-1}
   \int_0^Tu_q(T,S)w_{0,N}'(S)\ell_N(S)Z(S)\,\dd S. \label{eq:principal-green-attracting-LP}
\end{align}
On the repelling branch the mixed lift is the fixed point of
\begin{align}
 Z(T)&=U_C(T,0)b_s-\varepsilon^{-1}
   \int_0^TU_C(T,S)z_{0,N}'(S)\mathfrak p(S)\,\dd S,\notag\\
 \mathfrak p(T)&=u_q(T,T_r)b_u+\varepsilon^{-1}
   \int_T^{T_r}u_q(T,S)w_{0,N}'(S)\ell_N(S)Z(S)\,\dd S.
                                                               \label{eq:principal-green-repelling-LP}
\end{align}
The loop estimate \cref{eq:principal-green-loop} gives existence and
uniqueness.  It also shows directly that the terminal
\(\mathfrak p\)-row is uniformly transverse to the scalar unstable mode.

By linearity, split the repelling datum as
\((b_s,b_u)=(b_s,0)+(0,b_u)\).  Close the first fixed point in the norm
weighted by \(e^{c\Phi_s}\), and the second in the norm weighted by
\(e^{c\Phi_u}\).  The fast kernel preserves both weights after decreasing
\(r_{\rm out}\), because
\(|w_{0,N}'|\le C r_{\rm out}\) while its decay rate is bounded below by
\(D\gamma/(2\varepsilon)\); the same
\(O(r_{\rm out}^2)\) loop closes both fixed points.  Thus the stable datum
does not acquire a growing terminal scalar tail, and the terminal datum does
not acquire a growing initial fast tail.  The attracting integral equations
close similarly in the \(e^{c\Phi_a}\)-weighted norm.  Consequently the two
integral equations give
\cref{eq:principal-green-attracting-layer,eq:principal-green-repelling-layer}.
Successive time differentiation is triangular after multiplying the
\(j\)-th derivative by \(\varepsilon^j\).  Finally,
\(\dd T=\dd r/q_0(r)\), \(|q_0|\asymp1\), and
\(|w_{0,N}'(r)|\asymp|r|\) give
\cref{eq:principal-green-action-comparison}.

The trace identities follow at \(T=0\) and \(T=T_r\) in the two integral
systems.  Subtracting the slow particular solution from any solution with a
prescribed trace leaves the corresponding homogeneous lift.  This proves
\cref{eq:principal-green-lifts,eq:principal-green-augmented-inverse} and
uniqueness.
\end{proof}

\begin{proof}[Proof of \cref{prop:physical-history-phase-quotient}]
Because the rescaled fold-time RFDE is autonomous, differentiation of the
base solution in fold time gives a solution of the variational equation.  At the
history level this says
\begin{equation}
 \mathcal U_*(s,u)
 \left(\theta\mapsto v_*'(u+\theta),w_*'(u)\right)
 =\left(\theta\mapsto v_*'(s+\theta),w_*'(s)\right).
                                                               \label{eq:phase-quotient-time-tangent}
\end{equation}
Moreover,
\[
 \varrho_N
 \left(\theta\mapsto v_*'(s+\theta),w_*'(s)\right)
 =\pi_N^Tv_*'(s)=r_*'(s).
\]
Division by the nonzero endpoint speed proves both the normalization in
\cref{eq:physical-normalized-tangent} and
\cref{eq:physical-tangent-transport}.  It also shows uniqueness among scalar
multiples of the time tangent.

Since \(\varrho_N\tau_*(s)=1\), the operator \(Q_*(s)\) is a projection with
\begin{equation}
 \operatorname{Ran}Q_*(s)=\ker\varrho_N,
 \qquad
 \ker Q_*(s)=\operatorname{span}\{\tau_*(s)\}.      \label{eq:phase-quotient-range-kernel}
\end{equation}
For \(x\in\ker\varrho_N\), use the variational cocycle identity and
\cref{eq:physical-tangent-transport} to compute
\begin{align*}
 &Q_*(s)\mathcal U_*(s,t)Q_*(t)\mathcal U_*(t,u)x\\
 &\quad=Q_*(s)\mathcal U_*(s,u)x
 -Q_*(s)\mathcal U_*(s,t)\tau_*(t)
       \varrho_N\mathcal U_*(t,u)x\\
 &\quad=Q_*(s)\mathcal U_*(s,u)x,
\end{align*}
because the second term is a scalar multiple of \(Q_*(s)\tau_*(s)=0\).
This is \cref{eq:physical-quotient-cocycle}.  The range--kernel identities
identify the process with the quotient process, and the triangle inequality
gives \cref{eq:physical-quotient-bound}.  Finally, normalizing any nonzero
multiple of the time tangent by \(\varrho_N\tau=1\) recovers exactly
\cref{eq:physical-normalized-tangent}; hence the construction contains no
choice of tangent scale or frozen eigendirection.
\end{proof}

\begin{proof}[Proof of \cref{prop:exact-resource-gauge-quotient}]
The history derivative in
\cref{eq:exact-outer-history-tangent} gives
\(E_s^c=\partial_r\mathcal H_{V,w,q}(r_s)\).  Since the autonomous
time tangent is \(\delta q(r_s)E_s^c\),
\cref{eq:phase-quotient-time-tangent} yields
\[
 \mathcal U_*(s,\sigma)\{\delta q(r_\sigma)E_\sigma^c\}
 =\delta q(r_s)E_s^c,
\]
which proves \cref{eq:resource-gauge-center-cocycle}.

For reference, the full fold-time variational equation along the tracker is
\begin{align}
 \xi'(s)={}&\delta^{-1}\mathcal J_*(s)\xi(s)
 -\delta^{-1}\omega(s)\mathbf1\notag\\
 &+\delta K\sum_{k=0}^mL_{k,N}(\eta)
       \{\xi(s)-\xi(s-\theta_k)\},\notag\\
 \omega'(s)={}&\delta\pi_N^T\xi(s).                 \label{eq:resource-gauge-full-variation}
\end{align}
At the current endpoint,
\[
 \xi(s)=u(s)+V'(r_s)a(s),\qquad
 \omega(s)=w'(r_s)a(s).
\]
Differentiate \(a=\omega/w'(r_s)\), use \(r_s'=\delta q(r_s)\), and
differentiate the exact resource identity
\[
 q(r)w'(r)=r-\delta^2\nu
 \quad\Longrightarrow\quad
 q'(r)w'(r)+q(r)w''(r)=1=\pi_N^TV'(r).
\]
The full variation equation then gives
\begin{align*}
 a'
 &=\frac{\delta\pi_N^T\{u+V'a\}}{w'}
   -a\frac{\delta q w''}{w'}\\
 &=\frac{\delta}{w'}\pi_N^Tu+\delta q'a,
\end{align*}
which is \cref{eq:resource-gauge-a}.

For \(-\theta_m\le\theta<0\), histories translate, so
\[
 (\partial_s-\partial_\theta)\xi(s+\theta)=0.
\]
The definition of \(e_s\), or direct differentiation of
\(\delta q(r_s)e_s(\theta)=v_*'(s+\theta)\), gives
\begin{equation}
 (\partial_s-\partial_\theta)e_s(\theta)
 =-\delta q'(r_s)e_s(\theta).                       \label{eq:resource-gauge-e-transport}
\end{equation}
Apply \(\partial_s-\partial_\theta\) to
\(\psi_s=\xi_s-e_sa\) and use
\cref{eq:resource-gauge-a,eq:resource-gauge-e-transport}; the two
\(\delta q'a\) terms cancel and give
\cref{eq:resource-gauge-transport}.

The normalized center history is not itself a solution but obeys the
transport law \cref{eq:resource-gauge-center-cocycle}.  Equivalently, the
full variational operator applied to \(E_s^c\) equals
\(\partial_sE_s^c+\delta q'(r_s)E_s^c\).  Substitute
\[
 (\xi_s,\omega)
 =(\psi_s,0)+aE_s^c
\]
in \cref{eq:resource-gauge-full-variation}, subtract the derivative of
\(aE_s^c\), and use \cref{eq:resource-gauge-a}.  All terms proportional to
\(a\), including the delayed tangent terms, cancel.  The voltage endpoint
of the remaining rank-one term is
\[
 -V'(r_s)\frac{\delta}{w'(r_s)}\pi_N^Tu(s)
 =-\delta H_s\pi_N^Tu(s).
\]
This proves \cref{eq:resource-gauge-boundary}.  It also proves directly that
\(\psi_s\) is unchanged by adding a multiple of \(E_s^c\).

We next solve the nonstandard transport equation along characteristics.
For fixed \(s,\theta\), set
\[
 F(\sigma)=\psi_{s+\sigma}(\theta-\sigma),
 \qquad \theta\le\sigma\le0.
\]
Then \(F(\theta)=u(s+\theta)\), \(F(0)=\psi_s(\theta)\), and
\cref{eq:resource-gauge-transport} gives
\[
 F'(\sigma)
 =-\delta h_{s+\sigma}(\theta-\sigma)
       \pi_N^Tu(s+\sigma).
\]
Integration proves
\[
 \psi_s(\theta)=u_s(\theta)
 -\delta\int_\theta^0h_{s+\sigma}(\theta-\sigma)
        \pi_N^Tu_s(\sigma)\,\dd\sigma,
\]
which is \cref{eq:resource-gauge-history-conjugacy}.  In particular,
\[
 \psi_s(-\theta_k)
 =u(s-\theta_k)-(\mathcal K_su_s)(-\theta_k).
\]
Substitution in \cref{eq:resource-gauge-boundary} proves
\cref{eq:resource-gauge-standard-rfde}; the Volterra term therefore has the
displayed positive sign inside the delay bracket.  This derivation is along
characteristics on the delay-enlarged time interval, so no derivative of a
formally time-dependent conjugacy has been discarded.

It remains to prove the uniform bounds.  Under
\cref{eq:resource-gauge-collar-bounds}, the ratios of \(q\), the vector
\(V'\), and hence \(e_s\) are bounded by common constants, while
\[
 \|h_s(\theta)\|_N\le\frac{C}{\rho_\delta}.
\]
If \(\theta\le\sigma\le0\), then
\(e^{2\kappa_*(\theta-\sigma)}\le1\).  Therefore
\begin{align*}
 e^{2\kappa_*\theta}\|(\mathcal K_s\phi)(\theta)\|_N
 &\le\frac{C\delta}{\rho_\delta}
   \int_\theta^0
      e^{2\kappa_*(\theta-\sigma)}
      e^{2\kappa_*\sigma}\|\phi(\sigma)\|_N\,\dd\sigma\\
 &\le C\frac{\delta\theta_m}{\rho_\delta}
       \|\phi\|^\sharp_{\kappa_*}.
\end{align*}
The definition of \(\rho_\delta\) proves
\cref{eq:resource-gauge-volterra-bound}.  For small \(\delta\), the bound is
at most \(1/2\), and the Neumann series proves
\cref{eq:resource-gauge-volterra-inverse}.

Conversely, let \(u\) solve
\cref{eq:resource-gauge-standard-rfde} and define
\(\psi_s=(\Id-\mathcal K_s)u_s\).  Reversing the characteristic calculation
shows that \(\psi\) satisfies
\cref{eq:resource-gauge-transport}; reversing the endpoint substitution
shows that it satisfies \cref{eq:resource-gauge-boundary}.  Solve the scalar
linear equation \cref{eq:resource-gauge-a} from any prescribed value of
\(a\), and reconstruct
\[
 (\xi_s,\omega(s))=(\psi_s,0)+a(s)E_s^c.
\]
The cancellations already proved above, now read in reverse, show that this
pair solves the full variational equation.  Thus
\cref{eq:resource-gauge-history-conjugacy,eq:resource-gauge-history-inverse}
give the asserted two-sided quotient coordinate equivalence.

Finally, \(|\omega|/|w'(r_s)|=|a|\) and
\(\xi_s=\psi_s+e_sa\), with \(\|e_s\|^\sharp_{\kappa_*}\le C\).
The two triangle inequalities give the claimed uniform equivalence between
\cref{eq:resource-gauge-anisotropic-norm} and
\(\|\psi_s\|^\sharp_{\kappa_*}+|a|\).
\end{proof}

\subsection{Why unnormalized backward convergence does not select a history}

The next elementary RFDE example makes precise why an outer-history theorem
needs a normalized boundary coordinate.  It rules out a logical inference;
it remains compatible with a separately imposed global anchor for the
particular network field of the main article.

\begin{proposition}[Unnormalized backward convergence and loss of jet control]
\label[proposition]{prop:backward-asymptotic-nonselection}
Fix \(0<\delta_*<1\), let
\(\tau:(0,\delta_*)\to(0,\infty)\) be arbitrary, and write
\(\tau_\delta=\tau(\delta)\).  On \(\R^2\) use
\(|(y,z)|_1=|y|+|z|\), and put
\begin{equation*}
 X_\delta=C([-\tau_\delta,0];\R^2),\qquad
 X^1_\delta=C^1([-\tau_\delta,0];\R^2),
\end{equation*}
with
\begin{equation*}
 \|\phi\|_{X_\delta}=\sup_\vartheta|\phi(\vartheta)|_1,
 \qquad
 \|\phi\|_{X^1_\delta}=\|\phi\|_{X_\delta}
       +\sup_\vartheta|\phi'(\vartheta)|_1.
\end{equation*}
Consider the RFDE
\begin{equation}
 \dot x(t)=\mathcal F_\delta(x_t),
 \qquad \mathcal F_\delta(\phi)=\delta^{-2}\phi(0),           \label{eq:nonselection-rfde}
\end{equation}
whose compatible \(C^1\) histories satisfy
\begin{equation}
 \mathcal M^1_\delta
 =\{\phi\in X^1_\delta:\phi'(0)=\delta^{-2}\phi(0)\}.        \label{eq:nonselection-manifold}
\end{equation}

Even after fixing the common phase row
\begin{equation}
 \phi_y(0)=1,                                        \label{eq:nonselection-phase-row}
\end{equation}
bounded complete backward extension and convergence to zero in the
unnormalized history norm do not select a history.  Indeed, for every
\(A\in\R\),
\begin{equation}
 x_{\delta,A}(t)=e^{t/\delta^2}(1,A),\qquad t\le0,   \label{eq:nonselection-family}
\end{equation}
solves \cref{eq:nonselection-rfde} and satisfies
\begin{equation}
 \sup_{t\le0}\|x_{\delta,A,t}\|_{X_\delta}<\infty,
 \qquad
 \|x_{\delta,A,t}\|_{X_\delta}\longrightarrow0
       \quad(t\to-\infty).                          \label{eq:nonselection-asymptotics}
\end{equation}

Even superalgebraic closeness in \(X^1_\delta\), together with smooth
parameter dependence for each fixed \(\delta\), does not imply a tame
\(C^1_\nu C^2_\eta\) history jet.  Fix \(0<a<b\), put
\(\omega_\delta=e^{b/\delta^2}\), and, on
\(K=[-1/2,1/2]^2\), define
\begin{align}
 A_\delta(\nu,\eta)
 &=e^{-a/\delta^2}(1+\nu)
   \{\sin(\omega_\delta\eta)+1-\cos(\omega_\delta\eta)\},\notag\\
 \Phi_\delta(\nu,\eta)(\vartheta)
 &=e^{\vartheta/\delta^2}
       (1,A_\delta(\nu,\eta)),
 \qquad -\tau_\delta\le\vartheta\le0.             \label{eq:nonselection-jet-family}
\end{align}
The formula extends to a \(C^\infty\) map from \(\R^2\) to the closed
linear space \(\mathcal M^1_\delta\); restrict it to \(K\).  Every member
obeys \cref{eq:nonselection-phase-row}, extends to a solution satisfying
\cref{eq:nonselection-asymptotics}, and, relative to
\(\Phi_\delta^{\rm ref}(\vartheta)=e^{\vartheta/\delta^2}(1,0)\),
\begin{equation}
 \sup_K\|\Phi_\delta-\Phi_\delta^{\rm ref}\|_{X^1_\delta}
 \le\frac92(1+\delta^{-2})e^{-a/\delta^2}
 =O(\delta^q)\quad\text{for every }q>0.              \label{eq:nonselection-flat-size}
\end{equation}
At \((\nu,\eta)=(0,0)\), however,
\begin{align}
 \|\partial_\eta\Phi_\delta\|_{X_\delta}
 &=e^{(b-a)/\delta^2},&
 \|\partial_{\eta\eta}\Phi_\delta\|_{X_\delta}
 &=e^{(2b-a)/\delta^2},\notag\\
 \|\partial_\nu\partial_\eta\Phi_\delta\|_{X_\delta}
 &=e^{(b-a)/\delta^2},&
 \|\partial_\nu\partial_{\eta\eta}\Phi_\delta\|_{X_\delta}
 &=e^{(2b-a)/\delta^2}.                              \label{eq:nonselection-jet-blowup}
\end{align}
In particular, for arbitrary fixed \(C>0\), \(M,m,c\ge0\), and \(p>0\),
the ratio of each of the four displayed derivative norms to the envelope
\begin{equation}
 C\delta^{-M}\langle S_\delta\rangle^m e^{cS_\delta},
 \qquad S_\delta=\sqrt{2p\log(1/\delta)},            \label{eq:nonselection-tame-envelope}
\end{equation}
tends to \(+\infty\) as \(\delta\downarrow0\).  Since
\(\|\cdot\|_{X_\delta}\le\|\cdot\|_{X^1_\delta}\), the same conclusion
holds a fortiori for the strong-history derivative norms.
\end{proposition}

\begin{proof}
The functional \(\mathcal F_\delta\) is bounded linear on \(X_\delta\),
so \cref{eq:nonselection-rfde} is an RFDE, in fact in the nondelayed
subclass.  Direct differentiation proves \cref{eq:nonselection-family}.
Its history segment is
\begin{equation*}
 x_{\delta,A,t}(\vartheta)=e^{(t+\vartheta)/\delta^2}(1,A),
\end{equation*}
and hence
\begin{equation*}
 \|x_{\delta,A,t}\|_{X_\delta}=(1+|A|)e^{t/\delta^2},
 \qquad
 \|x_{\delta,A,t}\|_{X^1_\delta}
 =(1+\delta^{-2})(1+|A|)e^{t/\delta^2}.
\end{equation*}
This proves \cref{eq:nonselection-asymptotics}; the distinguished histories
\(x_{\delta,A,0}\) all satisfy the same phase row, while present evaluation
of the second component recovers the free coefficient \(A\).  Moreover,
\(\partial_\vartheta\phi(0)=\delta^{-2}\phi(0)\), so these histories and
the family below lie in \(\mathcal M^1_\delta\), which is the kernel of a
bounded linear map on \(X^1_\delta\) and hence is closed.

For the second component of \cref{eq:nonselection-jet-family}, relative to
the reference history, use
\(|\sin u|\le1\), \(|1-\cos u|\le2\), and
\(1+\nu\le3/2\).  This gives \cref{eq:nonselection-flat-size}, since
\(e^{-a/\delta^2}\) dominates every power of \(\delta\).  At \(\eta=0\),
direct differentiation of \(A_\delta\) gives coefficients
\begin{equation*}
 e^{(b-a)/\delta^2},\quad e^{(2b-a)/\delta^2},\quad
 e^{(b-a)/\delta^2},\quad e^{(2b-a)/\delta^2}
\end{equation*}
in the four derivative slots of
\cref{eq:nonselection-jet-blowup}.  The exponential profile has
\(X_\delta\)-norm one, proving the equalities.  Finally, the logarithm of
the ratio of the first quantity to
\cref{eq:nonselection-tame-envelope} is
\begin{equation*}
 \frac{b-a}{\delta^2}-M\log(1/\delta)
 -m\log\langle S_\delta\rangle-cS_\delta-\log C,
\end{equation*}
which tends to \(+\infty\).  The two derivative norms with exponent
\((2b-a)/\delta^2\) diverge still faster, proving all four assertions.
\end{proof}

Thus the words ``complete backward'', ``bounded'', and ``convergent in an
unnormalized history norm'', even with the displayed phase row, do not by
themselves select a complete history.  Nor may tame parameter jets be
inferred from fixed-\(\delta\) smoothness and superalgebraic \(C^0\) or
\(C^1\) closeness.  A valid construction must add a model-specific selection
or normalization and prove its uniform jets.
